%=========================
% Plantilla LaTeX Artículo Académico
% Compilar con: pdflatex + biber + pdflatex + pdflatex
%=========================
\documentclass[11pt,a4paper]{article}

%-------------------------
% Idioma y tipografía
%-------------------------
\usepackage[english]{babel}
\usepackage{csquotes}
\usepackage[T1]{fontenc}
\usepackage{lmodern}

%-------------------------
% Márgenes y espaciado
%-------------------------
\usepackage[a4paper,margin=2.5cm]{geometry}
\usepackage{setspace}
\onehalfspacing

%-------------------------
% Matemáticas y símbolos
%-------------------------
\usepackage{amsmath,amssymb,amsthm}

%-------------------------
% Tablas y figuras
%-------------------------
\usepackage{graphicx}
\usepackage{booktabs}
\usepackage{caption}
\usepackage{subcaption}
\usepackage{changepage}

%-------------------------
% Enlaces y PDF
%-------------------------
\usepackage{hyperref}
\hypersetup{
  colorlinks=true,
  linkcolor=blue,
  citecolor=blue,
  urlcolor=blue,
  pdftitle={The Continuum Hypothesis as a speech act},
  pdfauthor={Rodrigo Blanco Betes}
}

%-------------------------
% Listas, código y utilidades
%-------------------------
\usepackage{enumitem}
\usepackage{xcolor}
\usepackage{csquotes} % recomendado con biblatex

%-------------------------
% Bibliografía (biblatex + biber)
%-------------------------
\usepackage[
  backend=biber,
  style=apa,
  sortcites=true,
  url=true,
  doi=true
]{biblatex}
\DeclareLanguageMapping{english}{english-apa}
\addbibresource{refs.bib}

%-------------------------
% Metadatos del artículo
%-------------------------
\title{\textbf{Notes and references}}

\author{Rodrigo Blanco Betes}


\date{\today}

%-------------------------
% Entornos de teoremas (opcional)
%-------------------------
\newtheorem{theorem}{Theorem}
\newtheorem{lemma}{Lemma}

\begin{document}



\section{Set theory}

\begin{quote}
\small
Tarski's ``undefinability of truth'' shows that one cannot do this with proper class models. Naively, working in $ZFC$, it is trivial that each axiom $\varphi$ of $ZFC$ is true in $V$, so that $V \models ZFC$, so $\mathrm{Con}(ZFC)$. So, $ZFC \vdash \mathrm{Con}(ZFC)$. By Gödel's Second Incompleteness Theorem, there must be a flaw somewhere in this argument, but Tarski showed that one can't even define ``$\varphi$ is true in $V$'' as a property of the set $\varphi$, so that you can't even say ``$V \models ZFC$''. This is different from set models, where you can use Tarski's definition of truth to write the statement $\exists \mathcal{U}\ \mathcal{U} \models ZFC$ in set theory; it just isn't provable in $ZFC$ by the Incompleteness Theorem. \textcite[84]{Kunen2011}
\end{quote}

\begin{quote}
  \small
  With urelements in play, Zermelo investigates the various normal domains, models of ZF when the membership relation is restricted to them. In Zermelo’s words, a normal domain has a “width” given by its basis consisting of urelements, and a “height” given by its characteristic, the supremum of ordinal numbers represented in it. To modern eyes versed in pure set theory, i.e. set theory without urelements and with one universe, Zermelo’s width and height seem arcane, but as one reads on in the article, it becomes a crucial feature of the applicability of set theory. For Zermelo a mathematical context has basic subject matter and then various possible set-theoretic super-structures built on top for the application of set-theoretic ideas and constructions. \textcite[p. 392]{kanamori2010introduction}
\end{quote}



\printbibliography


\end{document}